\documentclass[10pt,a4paper]{amsart}
\usepackage[margin=19mm]{geometry}
\usepackage{amsmath,amssymb,mathtools}
\usepackage[scr=rsfs,cal=euler]{mathalfa}
\usepackage{xfrac}
\usepackage[unicode,hidelinks,bookmarks=false]{hyperref}
\newcommand{\ie}{i.e.\ }
\newcommand{\cf}{cf.\ }
\newcommand{\resp}{respectively\ }
\newcommand{\Subsection}{\subsection}
\providecommand{\nobreakdash}{\nobreak\hbox{-}}
\setlength{\emergencystretch}{2em}
\multlinegap0pt
\usepackage[shortlabels]{enumitem}
\usepackage{algorithm}\usepackage{algpseudocode}
\usepackage{amssymb}
\usepackage{mathtools}
\usepackage{xcolor}
\usepackage{tikz}
\newtheorem{assumption}{Assumption}
\newtheorem{theorem}{Theorem}
\newcommand{\N}{\mathbb{N}}
\newcommand{\abs}[1]{\left\lvert#1\right\rvert}
\renewcommand{\theassumption}{\Alph{assumption}}
\title{A priori adaptive numerical methods for estimating blow-up times of autonomous ODEs}
\author{H{\aa}kon Hoel, Johannes Vincent Meo}
\date{Source: arXiv:2603.12957v1 (CC BY 4.0)}
\begin{document}
\maketitle

\noindent\emph{Source and licence.} A priori adaptive numerical methods for estimating blow-up times of autonomous ODEs. Version arXiv:2603.12957v1 (CC BY 4.0), distributed by arXiv under the Creative Commons Attribution 4.0 International licence, http://creativecommons.org/licenses/by/4.0/, as recorded in the arXiv licence metadata for this exact version and shown on the abstract page https://arxiv.org/abs/2603.12957v1. Full licence notice: \texttt{licenses/arxiv-2603.12957-CC-BY-4.0.txt}.\par\smallskip

\noindent\textit{Editorial scope.} Counted unit: the article's theorem thm:error\_alg1 (source0.tex lines 216--220). The article's complete original proof of it is delivered with it (source0.tex lines 221--310, one \texttt{proof} environment of 90 lines). No mathematics is written, repaired or re-derived: the statement and the proof are the article's own lines, in the article's own notation, and no retained line is substituted or re-flowed.  Retained uncounted as the source-local context the proof needs: lines 129--148; lines 169--171; lines 184--187; lines 197--214.  Nothing was omitted from any retained span, so the package carries no omission record.  No retained line contains a citation, so the wrapper carries no bibliography and no bibliography entry.  No retained line contains a figure environment, an image-inclusion command or drawing code, so no image is delivered and the package declares no asset dependency.  Editorial formatting only, in addition: an AMS \texttt{amsart} preamble that restates the article's own declarations and macros used by the retained lines, namely the environments \texttt{assumption}, \texttt{theorem} on separate counters, together with the standard packages those lines need.  The retained environments print in this document as Assumption 1, Theorem 1 in order of appearance, whereas the article prints the same items with the numbering of its own full document.  The complete native source of the article as distributed in the arXiv e-print for this exact version is bundled unchanged as \texttt{sources/196306/source0.tex}.
We consider the autonomous one-dimensional ODE
\begin{equation}
\label{eq:ode}
\begin{split}
    x'(t) &= b(x(t)), \qquad t \in (0, \tau)\\
    x(0) &= x_0,
\end{split}
\end{equation}
for right-hand sides that satisfy the following assumption. 
\begin{assumption}
The function \(b:D = [x_0, \infty) \to (0,\infty)\) with \(x_0>0\) satisfies
\begin{enumerate}[label=\arabic*), ref=\theassumption.\arabic*]
    \label{asmp:b}
    \item \(b \in C^1(D)\) with \(b'\) positive and increasing on \(D\). \label{asmp:b.C1_bDerPos}
    \item There exist positive monotone decreasing functions \(F, G\in C^1((0,\infty))\) such that \(-\frac{d}{dx}G(b(x)) \leq\frac{1}{b(x)} \leq -\frac{d}{dx}F(b(x))\) for all \(x\in D\) with 
    \mbox{\(\underset{x\to\infty}{\lim}F(x)=0\)} and \(-\frac{1}{F^{-1}(\epsilon)G'(F^{-1}(\epsilon))}=\mathcal{O}(\epsilon^{\alpha})\) for some \(\alpha>-2\) for sufficiently small \(\epsilon>0\). \label{asmp:b.FandG}
    \item There exist some \(k>1\) such that \(\int_{x_0}^\infty\frac{\sqrt{b'(kx)}}{b(x)} \, dx <\infty\). \label{asmp:b.integrabilityCond}
    \item The integral \(\tau = \int_{x_0}^\infty \frac{1}{b(x)}dx\) is finite. \label{asmp:b.tauFinite}
\end{enumerate}
\end{assumption}

\begin{equation}\label{eq:tau_1D}
\tau_r(\tilde x_0, t_0) = t_0 + \int_{\min\{\tilde x_0, r\}}^r\frac{1}{b(x)}\,dx,
\end{equation}

\begin{equation}
\label{eq:xBarBndX}
    \bar{x}_{n+1} = \bar x_n + h_nb(\bar x_n) = x(t_n) + h_nb(x(t_n)) \leq x(t_n) + \int_{t_n}^{t_{n+1}}b(x(t))\,dt = x(t_{n+1}),
\end{equation}

\begin{algorithm}
    \caption{Estimating blow-up time of 1D ODE}
    \label{alg:1d}
    \begin{algorithmic}[1]
        \Require \(b,\, x_0,\, k,\, F^{-1},\,\epsilon > 0\).
        \State Set \(t = 0,\, N=0,\,\bar x = x_0\) and let \(r\) be such that \(b(r) = F^{-1}(\epsilon)\).
        \While{\(\bar x < r\)}
            \State \(h = \sqrt{\frac{\epsilon^2}{b'(\min\{k\bar x,\,r\})}}\).
            \State \(\bar x \leftarrow \bar x + b(\bar x)h\).
            \State \(t \leftarrow t + h\).
            \State \(N \leftarrow N+1\).
        \EndWhile
        \State \(\bar\tau = t\).
        \Ensure \(\bar\tau,\,N\).
    \end{algorithmic}
\end{algorithm}
As the number of floating point operations per iteration in Algorithm \ref{alg:1d} is bounded by a small fixed integer \(C\), we have that the total number of operations are \(\mathcal{O}(N)\). We will therefore define the computational cost to be the number of iterations \(N\).
The main result on the performance of our 1D numerical method is stated below. 

\begin{theorem}[Properties of Algorithm \ref{alg:1d}]
    \label{thm:error_alg1}
    Let Assumption~\ref{asmp:b} hold and let \(\bar\tau\) be calculated using Algorithm~\ref{alg:1d}.
    Then, for sufficiently small \(\epsilon > 0\), it holds that \(\abs{\bar\tau-\tau} = \mathcal{O}(\epsilon)\) and the computational cost of the method is \(\mathcal{O}\left(\epsilon^{-1}\right)\).
\end{theorem}

\begin{proof}
    For any \(r \geq x_0\) we have
    \begin{equation}
        \label{eq:errorSplit}
        \abs{\bar\tau-\tau} \leq \abs{\bar\tau-\tau_r} + \abs{\tau_r-\tau}.
    \end{equation}
    To prove our claim, it suffices to show both terms on the right-hand side are \(\mathcal{O}(\epsilon)\). 
    From Algorithm \ref{alg:1d} we have that \(r(\epsilon)\) is chosen such that \(b(r(\epsilon)) = F^{-1}(\epsilon)\), where we write \(r\) as a function of \(\epsilon\) to emphasize its dependence of the tolerance.
    For the second term, we then get
    \begin{equation}
        \label{eq:firstErrorTerm}
            \abs{\tau_{r(\epsilon)}-\tau} = \int_{r(\epsilon)}^\infty\frac{1}{b(x)}\,dx \stackrel{\ref{asmp:b.FandG}}{\leq} \int_{r(\epsilon)}^\infty-\frac{d}{dx}F(b(x)) \,dx = F(b(r(\epsilon)))=F(F^{-1}(\epsilon))=\epsilon.
    \end{equation}
    For later reference, from the lower bound of \(\frac{1}{b(x)}\) in Assumption \ref{asmp:b.FandG} we have that \(-\frac{d}{dx}G(b(x)) = -G'(b(x))b'(x) \leq \frac{1}{b(x)}\), thus by setting \(x = r(\epsilon)\) we get
    \begin{equation}
    \label{eq:boundB_der_r}
        b'(r(\epsilon)) \leq -\frac{1}{b(r(\epsilon))G'(b(r(\epsilon)))} = -\frac{1}{F^{-1}(\epsilon)G'(F^{-1}(\epsilon))} \stackrel{\ref{asmp:b.FandG}}{=} \mathcal{O}(\epsilon^\alpha),
    \end{equation}
    for some \(\alpha>-2\).

    For the first term of \eqref{eq:errorSplit}, we have
    \begingroup 
    \allowdisplaybreaks
    \begin{align}
    \label{eq:secondTermExp}
        \abs{\bar\tau-\tau_{r(\epsilon)}} &= \abs{\tau_{r(\epsilon)}(\bar{x}_N, t_N)-\tau_{r(\epsilon)}(\bar{x}_0, t_0)}\nonumber\\
        & \leq \sum_{n=1}^{N-1}\abs{\tau_{r(\epsilon)}(\bar{x}_n, t_n) - \tau_{r(\epsilon)}(\bar{x}_{n-1}, t_{n-1})} \nonumber \\
        &\hspace{6.5mm} + \abs{\tau_{r(\epsilon)}(\bar{x}_N, t_N) - \tau_{r(\epsilon)}(\bar{x}_{N-1}, t_{N-1})} \nonumber\\
        &\stackrel{\eqref{eq:tau_1D}}{\leq} \sum_{n=1}^{N-1}\abs{h_{n-1} - \int_{\bar x_{n-1}}^{\bar x_n}\frac{1}{b(x)}\, dx} + \abs{h_{N-1} - \int_{\bar x_{N-1}}^{r(\epsilon)}\frac{1}{b(x)}\, dx} \nonumber \\
        &\hspace{-1mm}\stackrel{\eqref{asmp:b.C1_bDerPos}}{\leq} \sum_{n=1}^{N-1}\abs{\left(\frac{1}{b(\bar x_{n-1})}-\frac{1}{b(\bar x_{n})}\right)b(\bar x_{n-1})h_{n-1}} + h_{N-1} \nonumber \\
        &\leq \sum_{n=1}^{N-1} b'(\bar x_n)h_{n-1}^2 + h_{N-1}.
    \end{align}
    \endgroup
    In the last inequality we take a first-order Taylor expansion of \(b(\bar x_n)\) around \(\bar x_{n-1}\) and use the monotonicity of $b'$ to $\max_{\theta \in [\bar x_{n-1}, \bar x_n]} |b'(\theta)| \le |b'(\bar x_n)|$.
    Moreover, one can show the the last line above is proportional to the sensitivity \(\partial_{x_0} \tau(r, x_0)\) times the local error in \(x\).
    
    The adaptive step length used in Algorithm \ref{alg:1d} gives \(h_{n-1} = \sqrt{\frac{\epsilon^2}{b'(\min\{k\bar x_{n-1},\,r(\epsilon)\})}}\), but for technical reasons, let us now introduce the following alternative step length:
    \begin{align}
        \label{eq:restrictions_h}
        g_{n-1} &= \min\Bigg\{ \underbrace{\sqrt{\frac{\epsilon^2}{b'(\min\{k\bar x_{n-1},\,r(\epsilon)\})}}}_{h_{n-1}}, \underbrace{\frac{(k-1)\bar x_{n-1}}{b(\bar x_{n-1})}}_{\tilde g_{n-1}} \Bigg\},
    \end{align}
    where \(k > 1\) is some value such that Assumption \ref{asmp:b.integrabilityCond} is satisfied. 
    Since \(g_{n-1} \leq \tilde g_{n-1}\), we have that
    \[\bar x_n = \bar x_{n-1} + g_{n-1}b(\bar x_{n-1}) \leq k\bar x_{n-1}.\] 
    Combining this with~\eqref{eq:secondTermExp} 
    and using that \(\bar x_{n} < r(\epsilon)\) for \(n < N\) yields
    \[\abs{\bar\tau-\tau_{r(\epsilon)}} \leq \sum_{n=1}^{N-1}b'(\min\{k\bar x_{n-1},\, r(\epsilon)\})g_{n-1}^2 + g_{N-1}\]
    Next, the other constraint on \(g_{n-1}\), \(h_{n-1}\), yields
    \begin{align}
        \label{eq:secondErrorTermFinal}
        \abs{\bar\tau-\tau_{r(\epsilon)}} \leq N\epsilon^2 + \mathcal{O}(\epsilon),
    \end{align}
    and it remains to bound \(N\). We first need to determine which of the constraints on \(g_{n-1}\) in \eqref{eq:restrictions_h} that is the most restrictive. For the quotient \(\frac{h_{n-1}}{\tilde g_{n-1}}\), we get
    \begin{align*}
        \frac{\sqrt{\frac{\epsilon^2}{b'(\min\{k\bar x_{n-1},\, r(\epsilon)\})}}}{\frac{(k-1)\bar x_{n-1}}{b(\bar x_{n-1})}} &= \frac{\epsilon b(\bar x_{n-1})}{(k-1)\bar x_{n-1}\sqrt{b'(\min\{k\bar x_{n-1},\, r(\epsilon)\})}} \\
        &\leq \frac{\epsilon c\bar x_{n-1} b'(\bar x_{n-1})}{(k-1)\bar x_{n-1}\sqrt{b'(\min\{k\bar x_{n-1},\, r(\epsilon)\})}} \\
        &\leq \frac{\epsilon c \sqrt{b'(r(\epsilon))}}{k-1} \\
        &\overset{\eqref{eq:boundB_der_r}}{=} \mathcal{O}(\epsilon^{1+\frac{\alpha}{2}}),
    \end{align*}
    where the transition to the second line uses that there exists a \(c>0\) such that \(b(x) = b(x_0) + b'(x)x \leq cxb'(x)\) for all \(x \in D\). 
    Thus for sufficiently small \(\epsilon\) we have that \(g_{n-1} = h_{n-1}\).

    Introducing the continuous extensions \(\bar x(t) = \bar x_{n-1}\) and \(h(t) = h_{n-1}\) for \(t \in [t_{n-1}, t_n)\), we obtain for sufficiently small \(\epsilon >0\) that
    \begingroup 
    \allowdisplaybreaks
    \begin{align}
    \label{eq:costEst1D}
        N &= \sum_{n=1}^{N}\frac{h_{n-1}}{h_{n-1}} \nonumber \\
        &\leq \int_0^{\bar\tau}\frac{1}{h(t)} \,dt \nonumber \\
          &\leq \frac{1}{\epsilon}\int_0^{\tau_{r(\epsilon)}} \sqrt{b'(k\bar x(t))} \,dt + \frac{1}{\epsilon}\int_{\tau_{r(\epsilon)}}^{\max\{\tau_{r(\epsilon)}, \bar\tau\}}\sqrt{b'(r(\epsilon))} \,dt \nonumber \\
          &\leq\frac{1}{\epsilon}\int_{x_0}^{\infty}\frac{\sqrt{b'(kx)}}{b(x)}\,dx + \frac{\abs{\bar\tau-\tau_{r(\epsilon)}}}{\epsilon}\sqrt{b'(r(\epsilon))} \nonumber \\
          &\stackrel{\ref{asmp:b.integrabilityCond}}{\le} \frac{C}{\epsilon} + \abs{\bar\tau-\tau_{r(\epsilon)}}\mathcal{O}(\epsilon^{\frac{\alpha}{2}-1}),
    \end{align}
    \endgroup
    where the transition from the second to the third line uses \eqref{eq:xBarBndX}, as this shows that \(\bar x(t) \leq x(t)\), allowing a change of variable in the left integral using \(dt = \frac{dx}{b(x)}\). 

    Inserting into~\eqref{eq:secondErrorTermFinal} gives
    \[\abs{\bar\tau-\tau_{r(\epsilon)}} \leq C\epsilon + \abs{\bar\tau-\tau_{r(\epsilon)}}\mathcal{O}(\epsilon^{1+\frac{\alpha}{2}}) \implies 
    \abs{\bar\tau-\tau_{r(\epsilon)}} 
    \leq \frac{C\epsilon}{(1-\mathcal{O}(\epsilon^{1+\frac{\alpha}{2}}))} 
    = C'\epsilon,\]
    as \(1+\frac{\alpha}{2}>0\), and thus
    \(N \stackrel{\eqref{eq:costEst1D}}{=} \mathcal{O}(\epsilon^{-1})\). 
    
    We then deduce from the above that 
    \[
    \abs{\bar\tau-\tau} \leq \abs{\bar\tau-\tau_{r(\epsilon)}} + \abs{\tau_{r(\epsilon)}-\tau} \leq C'\epsilon + \epsilon = \mathcal{O}(\epsilon).
    \]
    Furthermore, the computational cost is given by \(C''N \leq \frac{C'C''}{\epsilon} = \mathcal{O}(\epsilon^{-1})\), for some constant \(C''\in\N\).
\end{proof}

\end{document}
